\documentclass[11pt,a4paper]{article}

\usepackage[T1]{fontenc}
\usepackage[utf8]{inputenc}
\usepackage[english]{babel}
\usepackage{geometry}
\usepackage{hyperref}
\usepackage{longtable}
\usepackage{tabularx}
\usepackage{booktabs}
\usepackage{array}
\usepackage{enumitem}

\geometry{margin=2.5cm}
\hypersetup{
    colorlinks=true,
    linkcolor=blue,
    urlcolor=blue,
    pdftitle={SGEP --- Implementation Summary},
}

\title{SGEP --- Implementation Summary\\
\large Primary School Management System}
\author{Technical documentation generated from the SGEP repository}
\date{June 18, 2026}

\begin{document}

\maketitle
\tableofcontents
\newpage

%==============================================================================
\section{Overview}
%==============================================================================

\textbf{SGEP} (Système de Gestion d'École Primaire --- Primary School Management System) is a school management application tailored to the Cameroonian context. The repository combines a Django REST backend exposing the business API under \texttt{/api/v1/}, a Next.js~16 frontend (App Router), and Appwrite Cloud as a document database (without Django SQL ORM).

Active development is carried on branch \textbf{\texttt{feature/backend-02}} (present locally and on \texttt{origin}), which notably integrates Appwrite Cloud migration, attendance management, OpenAPI documentation, and the administrator dashboard.

%==============================================================================
\section{Technical Architecture}
%==============================================================================

\subsection{Stack}

\begin{table}[h]
\centering
\begin{tabular}{lll}
\toprule
\textbf{Layer} & \textbf{Technology} & \textbf{Role} \\
\midrule
Backend API & Django 5.2 LTS & Web framework \\
REST API & Django REST Framework 3.15+ & JSON endpoints \\
Persistence & Appwrite Cloud (Python SDK 7.x) & Document collections \\
Cache / broker & Redis 7 & Cache and Celery queue \\
Async tasks & Celery 5.4 + django-celery-beat & Exports, emails, report cards \\
API authentication & djangorestframework-simplejwt 5.x & JWT access / refresh \\
API documentation & drf-spectacular & OpenAPI schema, Swagger, ReDoc \\
PDF & WeasyPrint 62+ & Report cards, receipts \\
Excel & pandas + openpyxl & Tabular exports \\
Frontend & Next.js 16 (App Router) & User interface \\
Frontend auth & Appwrite SDK (session) + HttpOnly cookies & User login \\
\bottomrule
\end{tabular}
\end{table}

\subsection{Layered Pattern}

All business logic follows the mandatory flow documented in \texttt{AGENT.md}:

\begin{center}
\texttt{View (DRF)} $\rightarrow$ \texttt{Service} $\rightarrow$ \texttt{Repository} $\rightarrow$ \texttt{Appwrite SDK}
\end{center}

\begin{itemize}[noitemsep]
    \item DRF \textbf{views} validate input via serializers, apply permissions, and delegate to services.
    \item \textbf{Services} encapsulate business rules (calculations, orchestration, Celery dispatch).
    \item \textbf{Repositories} are the sole access point to Appwrite collections (\texttt{config/appwrite\_client.py}).
    \item No Django ORM models are used for business data (except \texttt{accounts.User} for JWT integration).
\end{itemize}

\subsection{Installed Django Apps}

Business applications registered in \texttt{INSTALLED\_APPS} are:

\begin{itemize}[noitemsep]
    \item \texttt{core} --- school, academic years, audit, Appwrite configuration
    \item \texttt{accounts} --- users, roles, JWT authentication, user dashboard
    \item \texttt{classes} --- classes and subjects
    \item \texttt{students} --- students, enrollments, promotions, exports
    \item \texttt{attendance} --- absences and tardiness
    \item \texttt{grades} --- grades and PDF report cards
    \item \texttt{finance} --- invoices, payments, financial dashboard
    \item \texttt{parents} --- parent portal and guardian accounts
    \item \texttt{notifications} --- emails and notification log (no dedicated REST URLs)
    \item \texttt{reports} --- tracking and download of async export jobs
\end{itemize}

\subsection{DRF Roles and Permissions}

Three business roles are defined in \texttt{accounts/permissions.py}:

\begin{table}[h]
\centering
\begin{tabular}{lll}
\toprule
\textbf{Role} & \textbf{Constant} & \textbf{Permission Class} \\
\midrule
Super Administrator & \texttt{superadmin} & \texttt{IsSuperAdmin} \\
Accountant & \texttt{comptable} & \texttt{IsComptable} (includes superadmin) \\
Parent & \texttt{parent} & \texttt{IsActiveParent} or \texttt{IsParentAny} \\
\bottomrule
\end{tabular}
\end{table}

\texttt{IsParentAny} allows access to invoices and receipts even when the parent account is suspended. \texttt{IsActiveParent} requires \texttt{account\_status = active} for academic data.

%==============================================================================
\section{Authentication --- Dual Flow}
%==============================================================================

\subsection{Appwrite Flow (Frontend)}

On login via \texttt{POST /api/auth/login} (Next.js Route Handler):

\begin{enumerate}[noitemsep]
    \item Call \texttt{loginWithAppwrite} --- creates an Appwrite session server-side.
    \item HttpOnly cookie \texttt{SESSION\_COOKIE} containing the Appwrite session secret.
    \item Non-HttpOnly CSRF cookie \texttt{psms\_csrf} for mutating requests.
\end{enumerate}

The endpoints \texttt{/api/auth/me}, \texttt{/api/auth/logout}, and session resolution rely on this Appwrite session.

\subsection{Django Simple JWT Flow (Business API Bridge)}

In parallel, the same login calls \texttt{loginWithDjango} toward \texttt{POST /api/v1/auth/login/}:

\begin{enumerate}[noitemsep]
    \item Email/password verification against the Appwrite \texttt{users} collection (\texttt{AuthService.login}).
    \item Issuance of an access/refresh JWT pair (Simple JWT) with claims \texttt{role}, \texttt{account\_status}, \texttt{email}, \texttt{student\_id}.
    \item HttpOnly cookies \texttt{DJANGO\_ACCESS\_COOKIE} and \texttt{DJANGO\_REFRESH\_COOKIE} stored by the frontend.
\end{enumerate}

Next.js Route Handlers (\texttt{proxyFetch}) inject the bearer JWT server-side; the browser never calls Django directly.

\subsection{Django API Request Authentication}

\texttt{REST\_FRAMEWORK.DEFAULT\_AUTHENTICATION\_CLASSES} uses \texttt{AppwriteJWTAuthentication}, which extends \texttt{JWTAuthentication} and loads the user from Appwrite via \texttt{UserRepository.get\_by\_id}.

Public auth endpoints (\texttt{authentication\_classes = []}): login and refresh. All other endpoints require a valid JWT unless stated otherwise.

%==============================================================================
\section{OpenAPI Documentation (Swagger)}
%==============================================================================

\begin{longtable}{@{}llp{7cm}@{}}
\toprule
\textbf{Path} & \textbf{Method} & \textbf{Description} \\
\midrule
\endhead
\texttt{/api/schema/} & GET & Raw OpenAPI schema (drf-spectacular) \\
\texttt{/api/docs/} & GET & Swagger UI interface \\
\texttt{/api/redoc/} & GET & ReDoc documentation \\
\texttt{/admin/} & --- & Django admin interface (outside API v1) \\
\bottomrule
\end{longtable}

Documented prefix for business routes: \texttt{/api/v1}.

%==============================================================================
\section{Backend API Routes (\texttt{/api/v1/})}
%==============================================================================

The tables below list routes declared in each module's \texttt{urls.py} files. The common prefix is \texttt{/api/v1/}.

\subsection{Module \texttt{accounts}}

\begin{longtable}{@{}lllp{5.5cm}@{}}
\toprule
\textbf{Route} & \textbf{Methods} & \textbf{Permission} & \textbf{Description} \\
\midrule
\endhead
\texttt{auth/login/} & POST & Public & Login, JWT issuance \\
\texttt{auth/refresh/} & POST & Public & Access token renewal \\
\texttt{auth/logout/} & POST & \texttt{IsAuthenticated} & Refresh token blacklist \\
\texttt{auth/change-password/} & POST & \texttt{IsAuthenticated} & Password change \\
\texttt{admin/dashboard/} & GET & \texttt{IsSuperAdmin} & User statistics (accounts, roles, recent) \\
\bottomrule
\end{longtable}

\textbf{Note:} \texttt{accounts.urls} is included first in \texttt{config/urls.py}. The \texttt{admin/dashboard/} route from \texttt{accounts} therefore takes precedence over the homonymous route in the \texttt{core} module (see Section~\ref{sec:admin-dashboard}).

\subsection{Module \texttt{core}}

\begin{longtable}{@{}lllp{5.5cm}@{}}
\toprule
\textbf{Route} & \textbf{Methods} & \textbf{Permission} & \textbf{Description} \\
\midrule
\endhead
\texttt{schools/} & GET, POST & \texttt{IsSuperAdmin} & List / create schools \\
\texttt{schools/<pk>/} & GET, PATCH & \texttt{IsSuperAdmin} & Detail / update \\
\texttt{academic-years/} & GET, POST & \texttt{IsSuperAdmin} & Academic years (filter \texttt{school\_id}) \\
\texttt{academic-years/<pk>/} & GET, PATCH & \texttt{IsSuperAdmin} & Detail / update \\
\texttt{admin/dashboard/} & GET & \texttt{IsSuperAdmin} & Enrollment dashboard (shadowed by \texttt{accounts}) \\
\bottomrule
\end{longtable}

\subsection{Module \texttt{classes}}

\begin{longtable}{@{}lllp{5.5cm}@{}}
\toprule
\textbf{Route} & \textbf{Methods} & \textbf{Permission} & \textbf{Description} \\
\midrule
\endhead
\texttt{classes/} & GET, POST & \texttt{IsSuperAdmin} & Classes (filters \texttt{academic\_year\_id}, \texttt{level\_id}) \\
\texttt{classes/<pk>/} & GET, PATCH & \texttt{IsSuperAdmin} & Detail / update \\
\texttt{subjects/} & GET, POST & \texttt{IsSuperAdmin} & Subjects (filter \texttt{class\_id}) \\
\texttt{subjects/<pk>/} & GET, PATCH & \texttt{IsSuperAdmin} & Detail / update \\
\bottomrule
\end{longtable}

\subsection{Module \texttt{students}}

\begin{longtable}{@{}lllp{5.5cm}@{}}
\toprule
\textbf{Route} & \textbf{Methods} & \textbf{Permission} & \textbf{Description} \\
\midrule
\endhead
\texttt{students/} & GET, POST & \texttt{IsSuperAdmin} & Paginated list / create student \\
\texttt{students/export/pdf/} & GET & \texttt{IsSuperAdmin} & Async PDF export (\texttt{202}, \texttt{job\_id}) \\
\texttt{students/export/excel/} & GET & \texttt{IsSuperAdmin} & Async Excel export (\texttt{202}, \texttt{job\_id}) \\
\texttt{students/<pk>/history/} & GET & \texttt{IsSuperAdmin} & Student history \\
\texttt{students/<pk>/enroll/} & POST & \texttt{IsSuperAdmin} & Class enrollment \\
\texttt{students/<pk>/promote/} & POST & \texttt{IsSuperAdmin} & Class promotion \\
\texttt{students/<pk>/} & GET, PATCH, DELETE & \texttt{IsSuperAdmin} & Detail / update / soft delete \\
\bottomrule
\end{longtable}

\subsection{Module \texttt{attendance}}

\begin{longtable}{@{}lllp{5.5cm}@{}}
\toprule
\textbf{Route} & \textbf{Methods} & \textbf{Permission} & \textbf{Description} \\
\midrule
\endhead
\texttt{attendance/} & GET, POST & \texttt{IsSuperAdmin} & List / record absence \\
\texttt{attendance/stats/} & GET & \texttt{IsSuperAdmin} & Attendance statistics \\
\texttt{attendance/export/} & GET & \texttt{IsSuperAdmin} & Async export (\texttt{format=pdf|excel}) \\
\texttt{attendance/<pk>/justify/} & POST & \texttt{IsSuperAdmin} & Absence justification \\
\texttt{attendance/<pk>/} & PUT & \texttt{IsSuperAdmin} & Update a record \\
\bottomrule
\end{longtable}

\subsection{Module \texttt{grades}}

\begin{longtable}{@{}lllp{5.5cm}@{}}
\toprule
\textbf{Route} & \textbf{Methods} & \textbf{Permission} & \textbf{Description} \\
\midrule
\endhead
\texttt{grades/} & GET, POST & \texttt{IsSuperAdmin} & List / create grade \\
\texttt{grades/bulk/} & POST & \texttt{IsSuperAdmin} & Bulk grade entry \\
\texttt{grades/<pk>/} & PUT & \texttt{IsSuperAdmin} & Update grade \\
\texttt{grades/export/results/} & GET & \texttt{IsSuperAdmin} & Async Excel results export \\
\texttt{report-cards/generate/} & POST & \texttt{IsSuperAdmin} & Async PDF report card generation \\
\texttt{report-cards/<pk>/status/} & GET & \texttt{IsSuperAdmin} & Report card job status \\
\texttt{report-cards/<pk>/download/} & GET & \texttt{IsSuperAdmin} or \texttt{IsActiveParent} & Report card download \\
\texttt{report-cards/<pk>/publish/} & POST & \texttt{IsSuperAdmin} & Publish report card (parent notification) \\
\bottomrule
\end{longtable}

\subsection{Module \texttt{finance}}

\begin{longtable}{@{}lllp{5.5cm}@{}}
\toprule
\textbf{Route} & \textbf{Methods} & \textbf{Permission} & \textbf{Description} \\
\midrule
\endhead
\texttt{finance/invoices/} & GET & \texttt{IsComptable} & Invoice list (filters \texttt{student\_id}, \texttt{status}) \\
\texttt{finance/invoices/generate/} & POST & \texttt{IsComptable} & Generate invoices by class \\
\texttt{finance/payments/} & POST & \texttt{IsComptable} & Record payment \\
\texttt{finance/payments/<pk>/receipt/} & GET & \texttt{IsComptable} & Receipt download \\
\texttt{finance/overdue/} & GET & \texttt{IsComptable} & Overdue invoices \\
\texttt{finance/dashboard/} & GET & \texttt{IsComptable} & Financial dashboard \\
\texttt{finance/export/excel/} & GET & \texttt{IsComptable} & Async Excel export \\
\bottomrule
\end{longtable}

\subsection{Module \texttt{parents} (Parent Portal)}

\begin{longtable}{@{}lllp{5.5cm}@{}}
\toprule
\textbf{Route} & \textbf{Methods} & \textbf{Permission} & \textbf{Description} \\
\midrule
\endhead
\texttt{parent/me/} & GET & \texttt{IsParentAny} & Parent profile (any status) \\
\texttt{parent/me/student/} & GET & \texttt{IsActiveParent} & Linked student record \\
\texttt{parent/me/grades/} & GET & \texttt{IsActiveParent} & Child's grades \\
\texttt{parent/me/attendance/} & GET & \texttt{IsActiveParent} & Child's absences \\
\texttt{parent/me/report-cards/} & GET & \texttt{IsActiveParent} & Report card list \\
\texttt{parent/me/report-cards/<pk>/download/} & GET & \texttt{IsActiveParent} & Report card download \\
\texttt{parent/me/invoices/} & GET & \texttt{IsParentAny} & Invoices (suspended account allowed) \\
\texttt{parent/me/invoices/<pk>/receipt/} & GET & \texttt{IsParentAny} & Payment receipt \\
\texttt{parent/me/messages/} & GET, POST & \texttt{IsActiveParent} & Admin $\leftrightarrow$ parent messaging \\
\bottomrule
\end{longtable}

\subsection{Module \texttt{reports}}

\begin{longtable}{@{}lllp{5.5cm}@{}}
\toprule
\textbf{Route} & \textbf{Methods} & \textbf{Permission} & \textbf{Description} \\
\midrule
\endhead
\texttt{reports/<pk>/status/} & GET & \texttt{IsComptable} if role is accountant/superadmin, else \texttt{IsSuperAdmin} & Export job tracking \\
\texttt{reports/<pk>/download/} & GET & Same & Generated file download \\
\bottomrule
\end{longtable}

\subsection{Async Export Pattern}
\label{sec:async-exports}

Export endpoints (students, attendance, results, finance) and report card generation follow the same pattern:

\begin{enumerate}[noitemsep]
    \item Creation of a \texttt{report\_jobs} document (status \texttt{pending}).
    \item Dispatch of a Celery task.
    \item HTTP \texttt{202 Accepted} response with \texttt{job\_id}.
    \item Polling via \texttt{GET /api/v1/reports/<job\_id>/status/}.
    \item Download via \texttt{GET /api/v1/reports/<job\_id>/download/}.
\end{enumerate}

%==============================================================================
\section{Administrator Dashboard}
\label{sec:admin-dashboard}
%==============================================================================

\subsection{Backend Endpoint}

\textbf{Effective route:} \texttt{GET /api/v1/admin/dashboard/} (\texttt{accounts} module, \texttt{IsSuperAdmin} permission).

The service \texttt{accounts.services.AdminDashboardService.get\_overview} returns:
\begin{itemize}[noitemsep]
    \item Counters: \texttt{total\_users}, \texttt{active\_users}, \texttt{suspended\_users}
    \item Breakdown by role: \texttt{superadmins}, \texttt{comptables}, \texttt{parents}
    \item \texttt{recent\_users} --- five most recently created accounts
\end{itemize}

A parallel implementation exists in \texttt{core.admin\_dashboard\_service.AdminDashboardService.build} (student, class, unpaid invoice, and recent activity statistics), registered on the same URL but unreachable while \texttt{accounts.urls} remains first.

\subsection{Frontend Bridge (branch \texttt{feature/backend-02})}

The Route Handler \texttt{GET /api/admin/dashboard} (Next.js) calls \texttt{proxyFetch("/admin/dashboard/")}, maps the Django response to the \texttt{AdminDashboardResponse} type, and feeds the \texttt{/admin} page.

%==============================================================================
\section{Frontend Routes}
%==============================================================================

\subsection{Route Handlers (\texttt{frontend/app/api/})}

\begin{longtable}{@{}llp{6.5cm}@{}}
\toprule
\textbf{Path} & \textbf{Methods} & \textbf{Role} \\
\midrule
\endhead
\texttt{/api/health} & GET & Availability probe (\texttt{\{status: "ok"\}}) \\
\texttt{/api/auth/login} & POST & Appwrite + Django JWT login, sets cookies \\
\texttt{/api/auth/logout} & POST & Logout (CSRF required), invalidates sessions \\
\texttt{/api/auth/me} & GET & Current user from Appwrite session \\
\texttt{/api/auth/role} & POST, DELETE & Role cookie (dev/transition) \\
\texttt{/api/admin/dashboard} & GET & Superadmin dashboard proxy \\
\bottomrule
\end{longtable}

\subsection{Application Pages (\texttt{frontend/app/})}

\begin{longtable}{@{}lp{8cm}@{}}
\toprule
\textbf{Path} & \textbf{Description} \\
\midrule
\endhead
\texttt{/} & Redirect to \texttt{/login} \\
\texttt{/login} & Login form (\texttt{(auth)/login}) \\
\texttt{/forgot-password} & Forgot password (Appwrite) \\
\texttt{/reset-password} & Password reset \\
\texttt{/admin} & Administrator dashboard (consumes \texttt{/api/admin/dashboard}) \\
\texttt{/admin/*} & Admin layout (sidebar navigation) \\
\texttt{/parent/*} & Parent portal layout (scaffolded) \\
\texttt{/accountant/*} & Accountant portal layout (scaffolded) \\
\bottomrule
\end{longtable}

\textbf{Frontend security:} Django JWTs and the Appwrite session are stored in HttpOnly cookies. The browser does not access the Django API directly; \texttt{proxyFetch} injects the bearer token server-side.

%==============================================================================
\section{Celery Tasks}
%==============================================================================

\subsection{Tasks by Module}

\begin{longtable}{@{}lp{9cm}@{}}
\toprule
\textbf{Task} & \textbf{Role} \\
\midrule
\endhead
\texttt{students.tasks.generate\_students\_export\_task} & Student PDF/Excel export \\
\texttt{attendance.tasks.generate\_attendance\_export\_task} & Attendance PDF/Excel export \\
\texttt{grades.tasks.generate\_class\_report\_cards\_task} & PDF report card generation by class \\
\texttt{reports.tasks.generate\_students\_excel\_task} & Student Excel export (reports module) \\
\texttt{reports.tasks.generate\_attendance\_excel\_task} & Attendance Excel export \\
\texttt{reports.tasks.generate\_results\_excel\_task} & Results Excel export \\
\texttt{reports.tasks.generate\_finance\_export\_task} & Finance Excel export \\
\texttt{finance.tasks.generate\_receipt\_task} & PDF/HTML receipt generation after payment \\
\texttt{finance.tasks.send\_overdue\_reminders\_task} & Overdue invoice reminders (email) \\
\texttt{parents.tasks.suspend\_inactive\_parents\_task} & Suspend parents not re-enrolled \\
\texttt{notifications.tasks.notify\_parent\_absence\_task} & Email on recorded absence \\
\texttt{notifications.tasks.notify\_bulletin\_published\_task} & Email when report card is available \\
\texttt{notifications.tasks.notify\_payment\_overdue\_task} & Overdue invoice email \\
\texttt{notifications.tasks.send\_parent\_credentials\_task} & Send parent credentials \\
\texttt{notifications.tasks.notify\_account\_suspended\_task} & Account suspended email \\
\texttt{notifications.tasks.notify\_account\_reactivated\_task} & Account reactivated email \\
\bottomrule
\end{longtable}

\subsection{Celery Beat Schedule}

Defined in \texttt{config/settings.py}:

\begin{itemize}[noitemsep]
    \item \texttt{suspend-inactive-parents-annual} --- September 1 at midnight (\texttt{parents.tasks.suspend\_inactive\_parents\_task})
    \item \texttt{send-overdue-reminders-daily} --- every day at 08:00 (\texttt{finance.tasks.send\_overdue\_reminders\_task})
\end{itemize}

Broker and results backend: Redis (\texttt{CELERY\_BROKER\_URL}, \texttt{CELERY\_RESULT\_BACKEND}).

%==============================================================================
\section{Appwrite Collections}
%==============================================================================

The management command \texttt{python manage.py setup\_appwrite} initializes database \texttt{APPWRITE\_DB\_ID} (default \texttt{sgep\_db}) with the following collections:

\begin{longtable}{@{}lp{8.5cm}@{}}
\toprule
\textbf{Collection ID} & \textbf{Domain} \\
\midrule
\endhead
\texttt{schools} & Schools \\
\texttt{academic\_years} & Academic years \\
\texttt{levels} & Cameroonian levels (SIL, CP, CE1, CM2, etc.) \\
\texttt{classes} & Classes by level and year \\
\texttt{students} & Student records \\
\texttt{student\_histories} & Student movement history \\
\texttt{parents} & Guardian data \\
\texttt{parent\_student} & Parent $\leftrightarrow$ student link \\
\texttt{users} & Application accounts (superadmin, accountant, parent) \\
\texttt{subjects} & Subjects and coefficients \\
\texttt{sequences} & Evaluation periods / sequences \\
\texttt{grades} & Grades \\
\texttt{attendance} & Absence records \\
\texttt{fee\_types} & School fee types \\
\texttt{invoices} & Invoices \\
\texttt{payments} & Payments \\
\texttt{report\_cards} & Generated report cards \\
\texttt{messages} & Admin/parent messaging \\
\texttt{notifications} & Sent notification log \\
\texttt{report\_jobs} & Async export jobs \\
\texttt{audit\_logs} & Audit log \\
\bottomrule
\end{longtable}

\texttt{AGENT.md} documents conceptual aliases (\texttt{guardians}, \texttt{enrollments}, \texttt{attendance\_records}); the effective identifiers created by \texttt{setup\_appwrite} are those in the table above.

%==============================================================================
\section{Tests}
%==============================================================================

\subsection{Covered Modules}

\begin{table}[h]
\centering
\begin{tabular}{lp{8cm}}
\toprule
\textbf{File / Module} & \textbf{Primary Coverage} \\
\midrule
\texttt{accounts/tests.py} & Login, refresh, logout, change-password JWT \\
\texttt{accounts/tests\_appwrite.py} & Appwrite user CRUD by role \\
\texttt{core/tests.py} & Audit, schools, repositories \\
\texttt{core/tests\_admin\_dashboard.py} & \texttt{core} dashboard service and endpoint \\
\texttt{classes/tests.py} & Classes and subjects \\
\texttt{students/tests.py} & Student CRUD, medical encryption, API \\
\texttt{attendance/tests.py} & Absences, stats, async export \\
\texttt{grades/tests.py} & Grades, bulk entry, report cards \\
\texttt{finance/tests.py} & Invoices, payments, dashboard \\
\texttt{parents/tests.py} & Parent portal, suspension, messages \\
\texttt{notifications/tests.py} & Email tasks (\texttt{dry\_run} mode) \\
\texttt{reports/tests.py} & Job status and download \\
\texttt{test\_auth\_login.py} & Root login test \\
\bottomrule
\end{tabular}
\end{table}

\subsection{Appwrite Integration Tests}

\texttt{core/tests\_appwrite\_integration.py} contains live tests against Appwrite Cloud, decorated with \texttt{@skip\_unless\_appwrite()}. They are skipped when environment variables contain placeholders.

Covered scenarios:
\begin{itemize}[noitemsep]
    \item Listing the \texttt{students} collection
    \item Full CRUD cycle on a test student
    \item Building the \texttt{core} dashboard under real conditions
\end{itemize}

%==============================================================================
\section{Branch \texttt{feature/backend-02}}
%==============================================================================

Branch \texttt{feature/backend-02} is the active SGEP backend working branch. Recent commits include:

\begin{itemize}[noitemsep]
    \item \texttt{develop} merge resolution and restoration of backend-02 work
    \item Attendance management features and API documentation
    \item Appwrite Cloud migration and student management enhancements
\end{itemize}

Notable deliverables on this branch include the \texttt{GET /api/v1/admin/dashboard/} endpoint (user statistics), its frontend proxy \texttt{GET /api/admin/dashboard}, the \texttt{/admin} page, and the full set of routes documented in this summary.

%==============================================================================
\section{Appendix --- Section Summary}
%==============================================================================

\begin{enumerate}
    \item Overview
    \item Technical architecture (stack, pattern, modules, roles)
    \item Authentication --- dual flow
    \item OpenAPI documentation (Swagger)
    \item Backend API routes (\texttt{/api/v1/})
    \item Administrator dashboard
    \item Frontend routes
    \item Celery tasks
    \item Appwrite collections
    \item Tests
    \item Branch \texttt{feature/backend-02}
\end{enumerate}

\end{document}
